\documentclass[11pt]{amsart}
\usepackage{amsmath,amssymb,amsthm,booktabs}
\usepackage[margin=1.15in]{geometry}
\usepackage{hyperref}
\newtheorem{theorem}{Theorem}
\newtheorem{proposition}[theorem]{Proposition}
\newtheorem{lemma}[theorem]{Lemma}
\newtheorem{corollary}[theorem]{Corollary}
\theoremstyle{definition}
\newtheorem{example}[theorem]{Example}
\newtheorem{remark}[theorem]{Remark}
\newcommand{\Prob}{\mathbb{P}}
\newcommand{\E}{\mathbb{E}}
\newcommand{\eps}{\varepsilon}
\newcommand{\cT}{\mathcal{T}}
\newcommand{\Cov}{\operatorname{Cov}}

\title[Harmonic measure on Galton--Watson trees at small bias]{The dimension of harmonic measure on
Galton--Watson trees increases at small bias}
\author{Yitzchak Shmalo}
\date{November 2026}

\begin{document}
\begin{abstract}
Let $T$ be a Galton--Watson tree without leaves whose offspring number $N$ is not constant and has finite mean
$m$, and let $d_\lambda$ be the Hausdorff dimension of the harmonic measure of the $\lambda$-biased random walk
on $T$, $0\le\lambda<m$. Lyons, Pemantle and Peres asked whether $d_\lambda$ increases with $\lambda$. We prove
that $\lambda\mapsto d_\lambda$ is continuously differentiable on $[0,k_0)$, where $k_0$ is the smallest value
taken by $N$, and that
\[
   \frac{d}{d\lambda}d_\lambda\Big|_{\lambda=0}=-\bigl(1-\E[1/N]\bigr)\Cov\bigl(\log N,1/N\bigr)>0 .
\]
Hence $d_\lambda$ is strictly increasing on an interval $[0,\eps]$, with $\eps>0$ depending on the offspring
law. No moment assumption beyond the finite mean is needed. The proof combines the formula of Lin and Rousselin
for $d_\lambda$ with a contraction estimate, uniform over trees, for the recursion satisfied by the escape
probabilities and by their derivatives in $\lambda$.
\end{abstract}
\maketitle

\section{Introduction}

Let $(p_k)_{k\ge1}$ be a probability distribution on the positive integers with $p_k<1$ for every $k$ and finite
mean $m=\sum_k kp_k$, so that $m>1$, and let $N$ denote a random variable with this distribution. Let $T$ be a
Galton--Watson tree with offspring distribution $(p_k)$, rooted at $e$. Every vertex of $T$ has at least one
child. We write $\nu(x)$ for the number of children of a vertex $x$ and
\[
   k_0=\min\{k\ge1:\ p_k>0\}
\]
for the smallest possible number of children.

For $\lambda\ge0$ the $\lambda$-biased random walk on $T$ starts at the root, moves from the root to a uniformly
chosen child, and moves from a vertex $x\ne e$ to its parent with probability $\lambda/(\lambda+\nu(x))$ and to
each of its children with probability $1/(\lambda+\nu(x))$. Lyons~\cite{Lyons90} proved that the walk is
transient exactly when $\lambda<m$. In that case the walk converges to a random ray of $T$; the law of this ray,
given $T$, is the $\lambda$-harmonic measure on the boundary $\partial T$ of $T$. The boundary carries the metric
$d(\xi,\eta)=e^{-|\xi\wedge\eta|}$, where $|\xi\wedge\eta|$ is the number of edges common to the rays $\xi$ and
$\eta$. Lyons, Pemantle and Peres~\cite{LPP95,LPP96} proved that the Hausdorff dimension of the
$\lambda$-harmonic measure is almost surely equal to a constant $d_\lambda$, and that $d_\lambda<\log m$, the
dimension of $\partial T$. For $\lambda=0$ the walk never backtracks, the harmonic measure is the visibility
measure of~\cite{LPP95}, and
\[
   d_0=\E[\log N].
\]

Lyons, Pemantle and Peres~\cite{LPP97} asked whether $d_\lambda$ is increasing in $\lambda$ on $[0,m)$, and
whether it is strictly increasing; this is Question~17.29 of~\cite{LyonsPeres}. Lin~\cite{Lin} and
Rousselin~\cite{Rousselin} found, independently, an explicit stationary measure for the environment seen along
the harmonic ray, and with it a formula for $d_\lambda$, recalled in Section~\ref{sec:proof}. Lin deduced that
$d_\lambda>d_0$ for every $\lambda\in(0,m)$, that $d_\lambda\to d_0$ as $\lambda\to0$ and that
$d_\lambda\to\log m$ as $\lambda\to m$. The comparison of $d_\lambda$ and $d_{\lambda'}$ for
$0<\lambda<\lambda'<m$ is left open in~\cite{Lin}, and the numerical computations of~\cite{Rousselin} support
monotonicity. To our knowledge these are the only results on the dependence of $d_\lambda$ on $\lambda$.
Lin~\cite{Lin} also observes that a neighbouring quantity, the mean number of children of the vertices along the
trajectory of the walk, is not monotone in $\lambda$.

The companion question for the speed of the walk, raised in~\cite{LPP96,LPP97}, is whether the speed on a
Galton--Watson tree without leaves is nonincreasing in $\lambda$ on $[0,m)$. It has been answered for small
$\lambda$. Ben Arous, Fribergh and Sidoravicius~\cite{BFS} proved monotonicity when the bias away from the root
is large enough. A\"{\i}d\'ekon~\cite{Aidekon} gave an explicit formula for the speed, and with it the range was
extended to $\lambda\le1/2$, a bound due to A\"{\i}d\'ekon and recorded in~\cite{SWX}, and then to
$\lambda\le k_0/\bigl(1+\sqrt{1-1/k_0}\bigr)$ when $k_0\ge2$ by Song, Wang and Xiang~\cite{SWX}. Mehrdad, Sen and
Zhu~\cite{MSZ} compare, for large bias, the speeds on two trees whose offspring distributions are stochastically
ordered. The regularity of the speed in $\lambda$ has been studied as well: Ben Arous, Hu, Olla and
Zeitouni~\cite{BHOZ} proved, under an exponential moment assumption, the Einstein relation, which gives the derivative of the speed at the critical
bias $\lambda=m$ in terms of the diffusivity of the critical walk, and Bowditch and Tokushige~\cite{BT} proved that
the speed is differentiable at every $\lambda$ for which the walk is ballistic and satisfies a central limit
theorem. All of these monotonicity results concern $\lambda<k_0$, which is also the range of
Theorem~\ref{thm:main} below.

\begin{theorem}\label{thm:main}
The function $\lambda\mapsto d_\lambda$ is continuously differentiable on $[0,k_0)$, with the one-sided
derivative at $0$, and
\begin{equation}\label{eq:slope}
   d_0'=\bigl(1-\E[1/N]\bigr)\Bigl(\E[1/N]\,\E[\log N]-\E\bigl[N^{-1}\log N\bigr]\Bigr)
       =-\bigl(1-\E[1/N]\bigr)\Cov\bigl(\log N,1/N\bigr) .
\end{equation}
This number is strictly positive.
\end{theorem}

\begin{corollary}\label{cor:mono}
There is $\eps>0$, depending on the offspring distribution, such that $d_\lambda<d_{\lambda'}$ whenever
$0\le\lambda<\lambda'\le\eps$.
\end{corollary}

Corollary~\ref{cor:mono} answers the question of~\cite{LPP97}, including strictness, on a neighbourhood of
$\lambda=0$. Since $1/N$ is bounded and $\E[\log N]\le m<\infty$, the right-hand side of~\eqref{eq:slope} is
finite under the standing assumptions; no second moment is used anywhere.

\begin{example}
If $N$ takes the values $1$ and $2$ with probability $1/2$ each, then $\E[1/N]=3/4$, $\E[\log N]=(\log 2)/2$ and
$\E[N^{-1}\log N]=(\log2)/4$, so $d_0'=(\log 2)/32$. If $N$ takes the values $2$ and $4$ with probability $1/2$
each, then $d_0'=5(\log2)/128$. For the geometric distribution $p_k=2^{-k}$ one has $\E[1/N]=\log 2$ and
$d_0'=0.05496$ to five decimal places.
\end{example}

\begin{remark}\label{rem:tilt}
In the formula of Lin and Rousselin, $d_\lambda$ is the mean of $\log(\lambda+C_\lambda(T))$ under a probability
measure on trees with an explicit density $W_\lambda$ with respect to the Galton--Watson measure; here
$C_\lambda(T)$ is the effective conductance of $T$. Both the function and the density depend on $\lambda$. At
$\lambda=0$ the derivative of $\log(\lambda+C_\lambda(T))$ is
$\bigl(1-\sum_{i\le\nu(e)}1/\nu(e_i)\bigr)/\nu(e)$, where $e_1,\dots,e_{\nu(e)}$ are the children of the root.
Its mean is zero, and it is negative with positive probability: this happens when the root has more than
$k_0$ children and each of them has $k_0$ children. The slope~\eqref{eq:slope} therefore comes entirely from the
change of the density, which to first order favours trees whose root has many children; see~\eqref{eq:W0}.
\end{remark}

The escape probability of the walk from a tree is the fixed point of a recursion over the subtrees of the
children of the root. Section~\ref{sec:escape} shows that, written for the logarithm of the escape probability
and restricted to trees in which every vertex has at least $k$ children, this recursion is a contraction in the
supremum norm with constant $\lambda/k$, uniformly in the tree, and that the same holds for the recursion
obtained by differentiating in $\lambda$. This gives continuous differentiability of the escape probability of
every such tree on $[0,k)$, with a bound on the derivative that does not depend on the tree, and allows
differentiation of the formula for $d_\lambda$ under the expectation. Section~\ref{sec:proof} contains the proof
of Theorem~\ref{thm:main}, and Section~\ref{sec:numerics} reports a numerical evaluation of $d_\lambda$ for four
offspring distributions.

\section{Escape probabilities}\label{sec:escape}

For an integer $k\ge1$ let $\cT_k$ be the set of rooted trees in which every vertex has finitely many and at
least $k$ children. For $T\in\cT_1$ let $T_*$ be the tree obtained by joining a new vertex $e_*$ to the root $e$
of $T$, so that $e_*$ is the parent of $e$. For $\lambda\ge0$ consider the walk on $T_*$ that starts at $e$ and
moves from every vertex $x\ne e_*$ to its parent with probability $\lambda/(\lambda+\nu(x))$ and to each child
with probability $1/(\lambda+\nu(x))$. For $j\ge0$ let
\begin{align*}
   b_j(T,\lambda)&=\Prob\bigl(\text{the walk reaches generation $j$ of $T$ before $e_*$}\bigr),\\
   \beta_\lambda(T)&=\Prob\bigl(\text{the walk never visits $e_*$}\bigr),
\end{align*}
so that $b_0=1$. If the root of $T$ has $n$ children, we write $T_1,\dots,T_n$ for the subtrees of $T$ rooted at
these children; they belong to $\cT_k$ when $T$ does.

\begin{lemma}\label{lem:rec}
Let $T\in\cT_1$, $\lambda\ge0$ and $j\ge1$, and let $S=\sum_{i=1}^n b_{j-1}(T_i,\lambda)$. Then
\begin{equation}\label{eq:rec}
   b_j(T,\lambda)=\frac{S}{\lambda+S},
\end{equation}
and $b_j(T,\lambda)$ decreases to $\beta_\lambda(T)$ as $j\to\infty$.
\end{lemma}

\begin{proof}
The first $j$ generations of $T$ form a finite set. For $\lambda>0$ the walk leaves this set almost surely,
because from each of its vertices the walk reaches $e_*$ in at most $j$ steps with probability bounded below;
for $\lambda=0$ the walk moves away from the root at every step. Hence, almost surely, the walk reaches $e_*$ or
generation $j$. Condition on the first step, which leads to $e_*$ with probability $\lambda/(\lambda+n)$ and to
the $i$th child with probability $1/(\lambda+n)$. From the $i$th child the walk reaches generation $j$ of $T$
before returning to $e$ with probability $b_{j-1}(T_i,\lambda)$, and otherwise it returns to $e$ and starts
afresh. Therefore $b=b_j(T,\lambda)$ satisfies
\[
   b=\frac1{\lambda+n}\sum_{i=1}^n\Bigl(b_{j-1}(T_i,\lambda)+\bigl(1-b_{j-1}(T_i,\lambda)\bigr)b\Bigr),
\]
that is $b(\lambda+S)=S$. Here $\lambda+S>0$: this is clear if $\lambda>0$, and if $\lambda=0$ then every
$b_{j-1}$ equals $1$ and $S=n$. So~\eqref{eq:rec} holds. The events in the definition of $b_j$ decrease with $j$, and
by the first sentence of the proof their intersection is, up to a null event, the event that the walk never
visits $e_*$.
\end{proof}

\begin{lemma}\label{lem:lower}
Let $k\ge1$, $T\in\cT_k$ and $0\le\lambda\le k$. Then $b_j(T,\lambda)\ge1-\lambda/k$ for every $j\ge0$, hence
$\beta_\lambda(T)\ge1-\lambda/k$. If $\lambda<k$, then
\begin{equation}\label{eq:beta}
   \beta_\lambda(T)=\frac{C_\lambda(T)}{\lambda+C_\lambda(T)},\qquad C_\lambda(T)=\sum_{i=1}^n\beta_\lambda(T_i) .
\end{equation}
\end{lemma}

\begin{proof}
The bound holds for $j=0$. If it holds for $j-1$ and all trees of $\cT_k$, then in~\eqref{eq:rec} we have
$S\ge n(1-\lambda/k)\ge k-\lambda$, and since $s\mapsto s/(\lambda+s)$ is nondecreasing,
$b_j(T,\lambda)\ge(k-\lambda)/k$. Letting $j\to\infty$ in~\eqref{eq:rec}, which involves finitely many subtrees
and a denominator bounded below by $k-\lambda>0$, gives~\eqref{eq:beta}.
\end{proof}

For $\lambda>0$ the quantity $C_\lambda(T)$ is the effective conductance between the root of $T$ and infinity
when an edge between generations $j$ and $j+1$ has conductance $\lambda^{-j}$.

\begin{proposition}\label{prop:C1}
Let $k\ge1$ and $0<r<k$, and put $c=1-r/k$ and $M=k/(k-r)^2$. For every $T\in\cT_k$ the function
$\lambda\mapsto\beta_\lambda(T)$ is continuously differentiable on $[0,r]$, and for all $\lambda\in[0,r]$,
\[
   c\le\beta_\lambda(T)\le1,\qquad \bigl|\partial_\lambda\beta_\lambda(T)\bigr|\le M,\qquad
   \partial_\lambda\beta_\lambda(T)\big|_{\lambda=0}=-\frac1{\nu(e)} .
\]
\end{proposition}

\begin{proof}
Throughout, $T$ ranges over $\cT_k$ and $\lambda$ over $[0,r]$, and $\|\cdot\|$ denotes the supremum over all
such pairs. Put $\rho=r/k<1$ and $x_j=\log b_j$. By Lemma~\ref{lem:lower}, $\log c\le x_j\le0$.

\emph{The recursion in logarithmic coordinates.} For $n\ge k$ and $x\in[\log c,0]^n$ let
$S(x)=\sum_ie^{x_i}$ and $H_\lambda(x)=\log S(x)-\log(\lambda+S(x))$, so that by~\eqref{eq:rec}
\[
   x_j(T,\lambda)=H_\lambda\bigl(x_{j-1}(T_1,\lambda),\dots,x_{j-1}(T_n,\lambda)\bigr).
\]
We have
\begin{equation}\label{eq:partials}
   \frac{\partial H_\lambda}{\partial x_i}=q\,w_i,\qquad q=\frac{\lambda}{\lambda+S},\quad w_i=\frac{e^{x_i}}{S},
   \qquad \frac{\partial H_\lambda}{\partial\lambda}=-\frac1{\lambda+S}.
\end{equation}
The $w_i$ are nonnegative and sum to $1$. Since $S\ge nc\ge k-r$, we have $q\le r/(r+k-r)=\rho$ and
$1/(\lambda+S)\le1/(k-r)$. The cube $[\log c,0]^n$ is convex, so $H_\lambda$ is $\rho$-Lipschitz on it for the
maximum norm. Consequently $\delta_j=\|x_j-x_{j-1}\|$ satisfies $\delta_j\le\rho\,\delta_{j-1}$ for $j\ge2$,
and $\delta_1=\sup\log((\lambda+n)/n)\le\log(1+\rho)\le\rho$. Thus $\delta_j\le\rho^j$.

\emph{The derivatives.} For fixed $T$ and $j$, the function $\lambda\mapsto b_j(T,\lambda)$ depends on the first
$j$ generations of $T$ only and, by~\eqref{eq:rec} and induction, is a rational function of $\lambda$ that is
positive on $[0,r]$. Let $y_j=\partial_\lambda x_j$, so that $y_0=0$. By the chain rule and~\eqref{eq:partials},
\begin{equation}\label{eq:yrec}
   y_j(T,\lambda)=-\frac1{\lambda+S_j}+q_j\sum_{i=1}^nw_{j,i}\,y_{j-1}(T_i,\lambda),
\end{equation}
where $S_j$, $q_j$ and $w_{j,i}$ are computed from the vector $(x_{j-1}(T_i,\lambda))_{i\le n}$. Hence
$\|y_j\|\le1/(k-r)+\rho\|y_{j-1}\|$, and by induction
\[
   \|y_j\|\le\frac1{(k-r)(1-\rho)}=M\qquad\text{for all }j.
\]
For $j\ge2$ subtract~\eqref{eq:yrec} for $j-1$ from~\eqref{eq:yrec} for $j$:
\[
   y_j-y_{j-1}=\Bigl(\frac1{\lambda+S_{j-1}}-\frac1{\lambda+S_j}\Bigr)
     +q_j\sum_iw_{j,i}\bigl(y_{j-1}-y_{j-2}\bigr)(T_i)
     +\sum_i\bigl(q_jw_{j,i}-q_{j-1}w_{j-1,i}\bigr)y_{j-2}(T_i).
\]
Since $|e^u-e^v|\le|u-v|$ for $u,v\le0$, we have $|b_{j-1}(T_i)-b_{j-2}(T_i)|\le\delta_{j-1}$ for every $i$ and
$|S_j-S_{j-1}|\le n\delta_{j-1}$. All of $S_j$, $S_{j-1}$ are at least $nc$. The first bracket is therefore at
most $n\delta_{j-1}/(nc)^2\le\delta_{j-1}/(kc^2)$ in absolute value. The second term is at most
$\rho\|y_{j-1}-y_{j-2}\|$. For the third, write $q\,w_i=\lambda b_i/(S(\lambda+S))$ with $b_i=e^{x_i}$, and let
$b_i'$, $S'$ be the corresponding quantities one step earlier. Then
\begin{align*}
  \sum_i\Bigl|\frac{b_i}{S(\lambda+S)}-\frac{b_i'}{S'(\lambda+S')}\Bigr|
  &\le\frac{\sum_i|b_i-b_i'|}{S(\lambda+S)}+S'\,\Bigl|\frac1{S(\lambda+S)}-\frac1{S'(\lambda+S')}\Bigr|\\
  &=\frac{\sum_i|b_i-b_i'|}{S(\lambda+S)}+\frac{|S-S'|\,(\lambda+S+S')}{S(\lambda+S)(\lambda+S')}
  \le\frac{n\delta_{j-1}}{(nc)^2}+\frac{2n\delta_{j-1}}{(nc)^2},
\end{align*}
where the last step uses $\lambda+S+S'\le(\lambda+S)+(\lambda+S')$. Multiplying by $\lambda\le r$, the third
term is at most $3r\delta_{j-1}M/(kc^2)$. Altogether, with $K=(1+3rM)/(kc^2)$ and
$\eta_j=\|y_j-y_{j-1}\|$,
\[
   \eta_j\le\rho\,\eta_{j-1}+K\delta_{j-1}\le\rho\,\eta_{j-1}+K\rho^{j-1}\qquad(j\ge2),
\]
and $\eta_1=\|y_1\|\le1/(k-r)$. By induction $\eta_j\le\rho^{j-1}\bigl(\eta_1+(j-1)K\bigr)$, so
$\sum_j\eta_j<\infty$.

\emph{Conclusion.} The functions $y_j$ converge uniformly on $\cT_k\times[0,r]$ to a function $y$ with
$\|y\|\le M$. For fixed $T$ each $y_j(T,\cdot)$ is continuous on $[0,r]$, hence so is $y(T,\cdot)$. By
Lemma~\ref{lem:rec}, $x_j(T,\lambda)\to x(T,\lambda)=\log\beta_\lambda(T)$. Passing to the limit in
$x_j(T,\lambda)=x_j(T,0)+\int_0^\lambda y_j(T,t)\,dt$, and noting that $x_j(T,0)=0$, we get
$x(T,\lambda)=\int_0^\lambda y(T,t)\,dt$. Thus $\lambda\mapsto\log\beta_\lambda(T)$ is continuously
differentiable on $[0,r]$ with derivative $y(T,\cdot)$, and
$|\partial_\lambda\beta_\lambda|=\beta_\lambda|y|\le M$. Finally, at $\lambda=0$ we have $b_{j-1}\equiv1$,
$S_j=n$ and $q_j=0$ in~\eqref{eq:yrec}, so $y_j(T,0)=-1/n$ for every $j\ge1$, and the same holds for
$y(T,0)=\partial_\lambda\beta_\lambda(T)|_{\lambda=0}$.
\end{proof}

\section{Proof of Theorem~\ref{thm:main}}\label{sec:proof}

Let $T$ and $T'$ be independent Galton--Watson trees with offspring distribution $(p_k)$. Almost surely both
belong to $\cT_{k_0}$, as do all the subtrees rooted at their vertices. For $0\le\lambda<k_0$ put
\[
   B_\lambda=\beta_\lambda(T'),\qquad C_\lambda=C_\lambda(T),\qquad
   W_\lambda=\frac{B_\lambda C_\lambda}{\lambda-1+B_\lambda+C_\lambda},\qquad L_\lambda=\log(\lambda+C_\lambda).
\]
For $0<\lambda<m$ the same definitions make sense, with $\beta_\lambda$ and $C_\lambda$ related as
in~\eqref{eq:beta}, and Lin~\cite[formulas (7) and (8) and the display that follows (9)]{Lin} proved that
\begin{equation}\label{eq:lin}
   d_\lambda=\frac{\E[W_\lambda L_\lambda]}{\E[W_\lambda]}\qquad(0<\lambda<m);
\end{equation}
see also~\cite[Theorem~4.1]{Rousselin}. The function $T\mapsto\E[W_\lambda\mid T]/\E[W_\lambda]$ is the density,
with respect to the Galton--Watson measure, of a probability measure on trees that is stationary for the
environment seen along the harmonic ray. At $\lambda=0$ we have $\beta_0\equiv1$, hence $C_0=\nu(e)$, $W_0=1$ and
$L_0=\log\nu(e)$, and the right-hand side of~\eqref{eq:lin} is $\E[\log N]=d_0$. So~\eqref{eq:lin} holds for
$\lambda=0$ as well.

\begin{proof}[Proof of Theorem~\ref{thm:main}]
Fix $r\in(0,k_0)$ and let $c=1-r/k_0$ and $M=k_0/(k_0-r)^2$. Write $n=\nu(e)$ for the number of children of the
root of $T$ and $D_\lambda=\lambda-1+B_\lambda+C_\lambda$. Primes denote derivatives in $\lambda$. By
Proposition~\ref{prop:C1}, applied to $T'$ and to $T_1,\dots,T_n$, the functions $B_\lambda$ and $C_\lambda$ are
continuously differentiable on $[0,r]$ and
\[
   c\le B_\lambda\le1,\qquad cn\le C_\lambda\le n,\qquad |B_\lambda'|\le M,\qquad |C_\lambda'|\le nM .
\]
By Lemma~\ref{lem:lower}, $B_\lambda\ge1-\lambda/k_0$ and $C_\lambda\ge n(1-\lambda/k_0)$, so
\[
   D_\lambda\ge\lambda\Bigl(1-\frac1{k_0}\Bigr)+n\Bigl(1-\frac{\lambda}{k_0}\Bigr)\ge cn,
   \qquad D_\lambda\le r+n .
\]
Therefore $W_\lambda$ and $L_\lambda$ are continuously differentiable on $[0,r]$, with
\[
   W_\lambda'=\frac{B_\lambda'C_\lambda+B_\lambda C_\lambda'}{D_\lambda}
     -\frac{B_\lambda C_\lambda(1+B_\lambda'+C_\lambda')}{D_\lambda^2},\qquad
   L_\lambda'=\frac{1+C_\lambda'}{\lambda+C_\lambda},
\]
and, using $n\ge1$,
\begin{gather*}
   \frac{c^2}{1+r}\le W_\lambda\le\frac1c,\qquad
   |W_\lambda'|\le\frac{2M}{c}+\frac{1+2M}{c^2},\\
   |L_\lambda|\le\log n+\kappa,\quad \kappa=\max\bigl(\log(1+r),\log(1/c)\bigr),\qquad
   |L_\lambda'|\le\frac{1+M}{c}.
\end{gather*}
Hence $|(W_\lambda L_\lambda)'|$ and $|W_\lambda'|$ are bounded on $[0,r]$ by a constant multiple of
$1+\log\nu(e)$, which is integrable because $\E[\log N]\le m$. By the mean value theorem and dominated
convergence, $\lambda\mapsto\E[W_\lambda L_\lambda]$ and $\lambda\mapsto\E[W_\lambda]$ are continuously
differentiable on $[0,r]$, with derivatives $\E[(W_\lambda L_\lambda)']$ and $\E[W_\lambda']$. Since
$\E[W_\lambda]\ge c^2/(1+r)$, the quotient~\eqref{eq:lin} is continuously differentiable on $[0,r]$. As
$r<k_0$ was arbitrary, this proves the first assertion, and
\begin{equation}\label{eq:quot}
   d_\lambda'=\frac{\E[W_\lambda'L_\lambda]+\E[W_\lambda L_\lambda']-d_\lambda\,\E[W_\lambda']}{\E[W_\lambda]}
   \qquad(0\le\lambda<k_0).
\end{equation}

It remains to evaluate~\eqref{eq:quot} at $\lambda=0$. Let $N=\nu(e)$ be the number of children of the root of
$T$, let $N_1,\dots,N_N$ be the numbers of children of these children, let $K$ be the number of children of the
root of $T'$, and put $A=\sum_{i\le N}1/N_i$. By Proposition~\ref{prop:C1},
\[
   B_0=1,\quad B_0'=-\frac1K,\qquad C_0=N,\quad C_0'=-A,\qquad D_0=N,\quad D_0'=1-\frac1K-A .
\]
Substituting,
\begin{equation}\label{eq:W0}
   W_0'=\frac{-N/K-A}{N}-\frac{N(1-1/K-A)}{N^2}=-\frac1K-\Bigl(1-\frac1K\Bigr)\frac1N,\qquad
   L_0'=\frac{1-A}{N};
\end{equation}
the terms containing $A$ cancel in $W_0'$. Let $a=\E[1/N]$ and $\ell=\E[\log N]$. The variable $K$ is independent
of $N$, and given $N$ the $N_i$ are independent copies of $N$, so $\E[A\mid N]=Na$. Hence
\[
   \E[L_0']=a-a=0,\qquad \E[W_0']=-2a+a^2,\qquad
   \E[W_0'\log N]=-a\ell-(1-a)\E\bigl[N^{-1}\log N\bigr].
\]
With $W_0=1$, $L_0=\log N$ and $d_0=\ell$, formula~\eqref{eq:quot} gives
\[
   d_0'=-a\ell-(1-a)\E\bigl[N^{-1}\log N\bigr]+(2a-a^2)\ell
       =(1-a)\Bigl(a\ell-\E\bigl[N^{-1}\log N\bigr]\Bigr),
\]
which is~\eqref{eq:slope}.

For the sign, let $N$ and $K$ be independent with law $(p_k)$. Then
\[
   -2\Cov\bigl(\log N,1/N\bigr)=\E\Bigl[\bigl(\log N-\log K\bigr)\Bigl(\frac1K-\frac1N\Bigr)\Bigr].
\]
The integrand is nonnegative, since the logarithm is increasing and $x\mapsto1/x$ is decreasing, and it is
strictly positive on the event $\{N\ne K\}$, which has positive probability because $N$ is not constant. Finally
$a<1$ because $\Prob(N\ge2)>0$. Hence $d_0'>0$.
\end{proof}

\begin{proof}[Proof of Corollary~\ref{cor:mono}]
By Theorem~\ref{thm:main} the derivative of $d_\lambda$ is continuous on $[0,k_0)$ and positive at $0$, so it is
positive on some interval $[0,\eps]$.
\end{proof}

\begin{remark}
Let $Q_\lambda$ be the probability measure with density $W_\lambda/\E[W_\lambda]$ with respect to the law of
$(T,T')$. Formula~\eqref{eq:quot} reads
\[
   d_\lambda'=\E_{Q_\lambda}\Bigl[\frac{1+C_\lambda'}{\lambda+C_\lambda}\Bigr]
     +\Cov_{Q_\lambda}\Bigl(\log(\lambda+C_\lambda),\frac{W_\lambda'}{W_\lambda}\Bigr)\qquad(0\le\lambda<k_0).
\]
The derivatives $B_\lambda'$ and $C_\lambda'$ are nonpositive, since the escape probabilities decrease in
$\lambda$, but this does not determine the sign of either term. At $\lambda=0$ the first term vanishes and the
second equals~\eqref{eq:slope}.
\end{remark}

\section{Numerical evaluation}\label{sec:numerics}

Formula~\eqref{eq:lin} can be evaluated by simulation of the recursion~\eqref{eq:beta}. A pool of $2\times10^5$
pairs, representing samples of $(\beta_\lambda(T),\partial_\lambda\beta_\lambda(T))$, is initialised at $(1,0)$
and updated $300$ times: every new pair is computed from~\eqref{eq:beta} and its derivative in $\lambda$, with an
independent offspring number and children drawn independently and uniformly from the current pool. The
expectations in~\eqref{eq:lin} and~\eqref{eq:quot} are then replaced by averages over the pool. Each run was
repeated with six seeds, and the same seeds were used for every $\lambda$.

The bias runs over the grid $\lambda/m\in\{0,\,0.02,\,0.05,\,0.1,\,0.2,\dots,0.9\}$, for the four offspring
distributions of Tables~\ref{tab:slope} and~\ref{tab:dim}.

\begin{table}[ht]\centering\small
\begin{tabular}{lccccc}\toprule
offspring law & $m$ & $d_0$ & $d_0'$ & pool average & quotient \\ \midrule
$p_1=p_2=\tfrac12$ & 1.500 & 0.34657 & 0.02166 & $0.02206\pm0.00041$ & $0.02228\pm0.00041$ \\
$p_2=p_4=\tfrac12$ & 3.000 & 1.03972 & 0.02708 & $0.02720\pm0.00019$ & $0.02718\pm0.00019$ \\
$p_1=p_2=p_3=\tfrac13$ & 2.000 & 0.59725 & 0.04954 & $0.04925\pm0.00020$ & $0.04954\pm0.00019$ \\
$p_k=2^{-k}$ & 2.000 & 0.50783 & 0.05496 & $0.05572\pm0.00062$ & $0.05677\pm0.00061$ \\
\bottomrule\end{tabular}
\caption{The slope at $\lambda=0$. The fourth column is~\eqref{eq:slope}, the fifth is the pool average of the
right-hand side of~\eqref{eq:quot} at $\lambda=0$, and the last is $(d_{\lambda}-d_0)/\lambda$ at
$\lambda=0.02\,m$; the last two are means over six seeds with their standard errors.}\label{tab:slope}
\end{table}

\begin{table}[ht]\centering\footnotesize\setlength{\tabcolsep}{4pt}
\begin{tabular}{lcccccccccc}\toprule
$\lambda/m$ & 0.0 & 0.1 & 0.2 & 0.3 & 0.4 & 0.5 & 0.6 & 0.7 & 0.8 & 0.9 \\ \midrule
$p_1=p_2=\tfrac12$ & 0.3466 & 0.3501 & 0.3540 & 0.3587 & 0.3643 & 0.3709 & 0.3787 & 0.3872 & 0.3954 & 0.4021 \\
$p_2=p_4=\tfrac12$ & 1.0394 & 1.0475 & 1.0556 & 1.0636 & 1.0713 & 1.0785 & 1.0850 & 1.0905 & 1.0948 & 1.0977 \\
$p_1=p_2=p_3=\tfrac13$ & 0.5978 & 0.6080 & 0.6190 & 0.6310 & 0.6436 & 0.6561 & 0.6676 & 0.6775 & 0.6854 & 0.6909 \\
$p_k=2^{-k}$ & 0.5080 & 0.5203 & 0.5356 & 0.5548 & 0.5786 & 0.6056 & 0.6324 & 0.6559 & 0.6747 & 0.6878 \\
\bottomrule\end{tabular}
\caption{Estimates of $d_\lambda$, means over six seeds. The standard errors, over the whole grid, are at most
$0.00050$.}\label{tab:dim}
\end{table}

In Table~\ref{tab:slope} the pool averages agree with~\eqref{eq:slope} to within one and a half standard errors,
and the difference quotients, taken at the positive bias $0.02\,m$, are within $0.002$ of it. In Table~\ref{tab:dim} the estimates increase along the grid and stay below
$\log m$, which equals $0.4055$, $1.0986$, $0.6931$ and $0.6931$ for the four distributions. The increase holds
for every one of the six seeds and every pair of consecutive grid points. This agrees with the computations
of~\cite{Rousselin} and with a positive answer to the question of~\cite{LPP97} on the whole interval $[0,m)$,
which Theorem~\ref{thm:main} leaves open for $\lambda$ away from $0$.

\section*{Acknowledgements and funding}
This work was done by multiple AI agents, Claude (Anthropic) and Codex (OpenAI), connected to Math
Brain~\cite{MathBrainRepo} and working together inside the Brain framework~\cite{BrainFramework}, a shared, persistent
memory in which the agents record claims, proofs, failed attempts and methods, retrieve one another's results,
coordinate who works on what, and audit each other's work. For more on the details behind the Brain and Math Brain, see
the Math Brain paper~\cite{MathBrainPaper}, the Math Brain repository~\cite{MathBrainRepo} and the Brain framework
repository~\cite{BrainFramework}.

The research presented in this paper was supported by the European Research Council (ERC) under the European Union's
Horizon Europe research and innovation programme (grant agreement No.~101041711), by the Simons Foundation as part of
the Collaboration on the Mathematical and Scientific Foundations of Deep Learning, by Heights Labs, by Convex Nexus
Capital, by the Israel Science Foundation (grant number 2258/19), and by the Israel Science Foundation (ISF Grant
4101/25). Computational resources were provided in part by Google Cloud through its education credits program (credit
530550103).

\section*{Declaration of competing interest}
The author holds equity in Heights Labs and in Convex Nexus Capital and has a management role at Convex Nexus Capital.
These entities are included in the funding declaration above.

\section*{Declaration of generative AI and AI-assisted technologies}
The paper comes out of Math Brain, an ongoing project on the use of AI systems for research mathematics, and its
results were developed with the Brain~\cite{BrainFramework} and Math Brain~\cite{MathBrainRepo}. In October 2026
AI agents working from the project's list of open problems in probability obtained Theorem~\ref{thm:main} and
proofs of it, and a further agent checked the arguments. In November 2026 a
Claude Fable agent (Anthropic) checked the proof again, wrote the simulation program and prepared the manuscript.
These checks are not peer review. The author directed the work and is responsible for the content of the paper and
for decisions concerning its dissemination.

\section*{Code and data availability}
Programs and outputs:\\
\url{https://github.com/yspennstate/harmonic-dimension-small-bias}.

The paper is part of the Math Brain project, whose public repository is
\url{https://github.com/yspennstate/math-brain-public}~\cite{MathBrainRepo}.

\begin{thebibliography}{99}
\bibitem{Aidekon} E.~A\"{\i}d\'ekon, Speed of the biased random walk on a Galton--Watson tree,
\emph{Probab. Theory Related Fields} \textbf{159} (2014), 597--617.
\bibitem{BFS} G.~Ben Arous, A.~Fribergh and V.~Sidoravicius, Lyons--Pemantle--Peres monotonicity problem for high
biases, \emph{Comm. Pure Appl. Math.} \textbf{67} (2014), 519--530.
\bibitem{BHOZ} G.~Ben Arous, Y.~Hu, S.~Olla and O.~Zeitouni, Einstein relation for biased random walk on
Galton--Watson trees, \emph{Ann. Inst. Henri Poincar\'e Probab. Stat.} \textbf{49} (2013), 698--721.
\bibitem{BT} A.~Bowditch and Y.~Tokushige, Differentiability of the speed of biased random walks on
Galton--Watson trees, \emph{ALEA Lat. Am. J. Probab. Math. Stat.} \textbf{17} (2020), 609--642.
\bibitem{Lin} S.~Lin, Harmonic measure for biased random walk in a supercritical Galton--Watson tree,
\emph{Bernoulli} \textbf{25} (2019), 3652--3672.
\bibitem{Lyons90} R.~Lyons, Random walks and percolation on trees, \emph{Ann. Probab.} \textbf{18} (1990),
931--958.
\bibitem{LPP95} R.~Lyons, R.~Pemantle and Y.~Peres, Ergodic theory on Galton--Watson trees: speed of random walk
and dimension of harmonic measure, \emph{Ergodic Theory Dynam. Systems} \textbf{15} (1995), 593--619.
\bibitem{LPP96} R.~Lyons, R.~Pemantle and Y.~Peres, Biased random walks on Galton--Watson trees,
\emph{Probab. Theory Related Fields} \textbf{106} (1996), 249--264.
\bibitem{LPP97} R.~Lyons, R.~Pemantle and Y.~Peres, Unsolved problems concerning random walks on trees, in
\emph{Classical and Modern Branching Processes}, IMA Vol. Math. Appl. \textbf{84}, Springer, 1997, 223--237.
\bibitem{LyonsPeres} R.~Lyons and Y.~Peres, \emph{Probability on Trees and Networks}, Cambridge University Press,
2016. Question numbers are those of the version on the first author's web page.
\bibitem{MSZ} B.~Mehrdad, S.~Sen and L.~Zhu, The speed of a biased walk on a Galton--Watson tree without leaves is
monotonic with respect to progeny distributions for high values of bias, \emph{Ann. Inst. Henri Poincar\'e
Probab. Stat.} \textbf{51} (2015), 304--318.
\bibitem{Rousselin} P.~Rousselin, Invariant measures, Hausdorff dimension and dimension drop of some harmonic
measures on Galton--Watson trees, \emph{Electron. J. Probab.} \textbf{23} (2018), paper no.~46, 1--31.
\bibitem{SWX} H.~Song, L.~Wang and K.~Xiang, The speed of a biased walk on a Galton--Watson tree without leaves is
monotonic for low values of bias, \emph{J. Appl. Probab.} \textbf{62} (2025), 1044--1052.
\bibitem{MathBrainRepo} Y.~Shmalo, \emph{Math Brain: public snapshot, code and tests}, 2026,
\url{https://github.com/yspennstate/math-brain-public}.
\bibitem{MathBrainPaper} Y.~Shmalo, \emph{Math Brain: shared, auditable memory for multi-agent mathematical research},
October 2026, \url{https://github.com/yspennstate/math-brain-public/tree/main/papers/math-brain-2026-10-24}.
\bibitem{BrainFramework} Y.~Shmalo, \emph{Brain Framework}, public software repository, 2026,
\url{https://github.com/yspennstate/brain-framework-public}.
\end{thebibliography}
\end{document}
